\documentclass[10pt,a4paper]{article}

% =========================================================
% One-Page Print Appendix Template
% Compact Editorial Style, No Colors
% =========================================================

\usepackage[a4paper,margin=1.02cm]{geometry}
\usepackage[T1]{fontenc}
\usepackage{lmodern}
\usepackage{microtype}
\usepackage{multicol}
\usepackage{amsmath,amssymb}
\usepackage{booktabs}
\usepackage{array}
\usepackage{enumitem}
\usepackage{titlesec}
\usepackage{graphicx}

% ---------------------------------------------------------
% Page setup
% ---------------------------------------------------------
\pagestyle{empty}
\setlength{\parindent}{0pt}
\setlength{\parskip}{2.2pt}
\setlength{\columnsep}{0.58cm}
\setlength{\columnseprule}{0.25pt}
\setlength{\fboxsep}{4pt}
\setlength{\abovedisplayskip}{3pt}
\setlength{\belowdisplayskip}{3pt}
\setlength{\abovedisplayshortskip}{2pt}
\setlength{\belowdisplayshortskip}{2pt}

\setlist[itemize]{
    leftmargin=1.05em,
    itemsep=0pt,
    topsep=1pt,
    parsep=0pt
}

\renewcommand{\thesection}{\Alph{section}}
\renewcommand{\thesubsection}{\thesection.\arabic{subsection}}

% ---------------------------------------------------------
% Compact section titles
% ---------------------------------------------------------
\titleformat{\section}
  {\large\bfseries}
  {\fbox{\strut APP.\,\thesection}}
  {0.45em}
  {}
  [\vspace{-3pt}\rule{\linewidth}{0.45pt}]

\titleformat{\subsection}
  {\normalsize\bfseries}
  {\thesubsection}
  {0.4em}
  {}

\titlespacing*{\section}{0pt}{5pt}{4pt}
\titlespacing*{\subsection}{0pt}{3pt}{1pt}

% ---------------------------------------------------------
% Print-friendly boxed note
% ---------------------------------------------------------
\newcommand{\printbox}[2]{%
\vspace{2pt}
\noindent\fbox{%
\begin{minipage}{\dimexpr\linewidth-2\fboxsep-2\fboxrule\relax}
\textbf{#1}\\[-1pt]
{\small #2}
\end{minipage}
}
\vspace{2pt}
}

% ---------------------------------------------------------
% Clear objective block
% ---------------------------------------------------------
\newcommand{\objectivebox}[1]{%
\vspace{4pt}
\noindent\fbox{%
\begin{minipage}{\dimexpr\linewidth-2\fboxsep-2\fboxrule\relax}
{\normalsize\bfseries OBJECTIVE}\\[-1pt]
{\small #1}
\end{minipage}
}
\vspace{4pt}
}

% ---------------------------------------------------------
% Header
% ---------------------------------------------------------
\newcommand{\AppendixTitle}[2]{%
\noindent
\begin{minipage}[c]{0.07\textwidth}
\centering
\rotatebox{90}{\large\bfseries APPENDIX}
\end{minipage}
\hfill
\begin{minipage}[c]{0.90\textwidth}
{\Huge\bfseries APPENDICES}\\[-2pt]
{\Large\scshape #1}\\[3pt]
\rule{\linewidth}{0.8pt}\\[-2pt]
{\small #2}
\end{minipage}

\objectivebox{This appendix is a compact technical reference for the material that supports the main thesis: pruning-space definitions, MAC accounting, fine-tuning objectives, implementation choices, and reporting rules. It is written to make the experimental setup readable at a glance, without interrupting the main discussion.}
}

% =========================================================
% Document
% =========================================================

\begin{document}

\AppendixTitle
{Compact Supplementary Material}
{Additional derivations, implementation details, experimental settings, and result-reporting conventions.}

\begin{multicols}{2}
\small

\section{Pruning Search Space}

\printbox{Allocation Vector}{
For a Transformer with \(L\) layers, \(H\) attention heads, and \(F\) feed-forward channels, a structured pruning candidate is encoded as
\[
s=(h_1,f_1,h_2,f_2,\ldots,h_L,f_L).
\]
Each pair \((h_\ell,f_\ell)\) specifies the retained structure in layer \(\ell\).
}

The global budget constraint is written as
\[
\operatorname{MAC}(s)\leq \Omega ,
\]
where \(\Omega\) is the target computational budget.

\subsection{Attention Heads}
For sequence length \(n\), hidden size \(d\), and head dimension \(d_h\), the approximate MAC cost of one retained head is
\[
3ndd_h + 2n^2d_h + nd_hd .
\]

\subsection{Feed-Forward Channels}
For one retained FFN channel, the approximate contribution is
\[
2nd,
\]
corresponding to the two linear projections in the feed-forward block.

\section{Fine-Tuning Objective}

\printbox{Training Rule}{
The searched structure is physically instantiated before fine-tuning. The teacher model remains fixed, while the compressed student is optimized under the selected architecture.
}

The total objective combines supervised loss and distillation loss:
\[
\mathcal{L}
=
\lambda_{\mathrm{ce}}\mathcal{L}_{\mathrm{ce}}
+
\lambda_{\mathrm{kd}}\mathcal{L}_{\mathrm{kd}} .
\]

A softened KL distillation term is written as
\[
\mathcal{L}_{\mathrm{kd}}
=
T^2
\operatorname{KL}
\left(
\sigma(z_t/T)
\;\|\;
\sigma(z_s/T)
\right),
\]
where \(z_t\) and \(z_s\) are teacher and student logits.

\section{Experimental Notes}

\printbox{Reporting Principle}{
When the main claim concerns computational efficiency, accuracy should be compared at equal MAC budget. Parameter count alone does not fully describe inference cost.
}

\subsection{Compact Result Table}

\begin{center}
\renewcommand{\arraystretch}{1.08}
\begin{tabular}{lccc}
\toprule
Task & MACs & Params & Score \\
\midrule
SST-2 & \(40\%\) & 35M  & 93.7 \\
QQP   & \(40\%\) & 35M  & 91.3 \\
MNLI  & \(8\%\)  & 3.5M & 81.8 \\
\bottomrule
\end{tabular}
\end{center}

\section{Implementation Details}

A candidate may be stored as two vectors:
\[
h=(h_1,\ldots,h_L), \qquad
f=(f_1,\ldots,f_L).
\]
This representation makes mutation, crossover, and budget repair straightforward.

\subsection{Budget Repair}
After mutation, an invalid candidate is repaired by removing additional units until
\[
\operatorname{MAC}(s)\leq\Omega .
\]

\printbox{Physical Pruning}{
Masks are converted into an actual compact model by removing selected heads and FFN channels. Fine-tuning is then applied directly to the reduced architecture.
}

\section{Checklist}

\begin{itemize}
    \item Report the exact MAC budget.
    \item State whether embeddings are counted.
    \item Use the same sequence length across comparisons.
    \item Separate search cost from fine-tuning cost.
    \item Compare against baselines at matched budgets.
\end{itemize}

\end{multicols}

\end{document}
